\chapter*{Before we start}
\addcontentsline{toc}{chapter}{Before we start}

\nodwyn{} is an agent that works in your folders: it reads, writes, runs your
tests and talks to a language model. What makes it different from the other
agents you may have tried is a promise it keeps no matter what the model says:
\textbf{nothing happens to your computer without passing one checkpoint}, and
everything that does happen is written down first. The model proposes; you
nod. That is where the name comes from.

The second thing that makes it different is that almost everything above that
checkpoint is written in a small, friendly language, \textbf{Scheme}, and you
can change it. The loop that asks the model and runs its tools, the prompts
it sends, the planner, the context compaction, the sub-agents, the habits
that make it good at fixing bugs: all of it is Scheme you can read, wrap,
replace or delete, from a file in your project.

This book takes you from ``what is this?'' to ``I rewired how it works in my
repository this morning.'' It is organised as a walk:

\begin{description}[leftmargin=1.2em,labelindent=0pt]
  \item[Part I, Meeting nodwyn] what it is, how to start it, and the
    different ways to reach it: the app, the browser, the terminal, your
    phone.
  \item[Part II, Scheme, the Fun Part] the language, from nothing, with
    puzzles. You do not need to know Lisp. You will by the end.
  \item[Part III, Scheme Meets the World] how Scheme touches files and
    programs through \emph{primitives}, why every one of them passes the
    \emph{guard}, and how to give the model tools of your own.
  \item[Part IV, Steering the Loop] the agent loop itself: turns,
    strategies, contexts, sub-agents, plans, runs, and the working habits
    that ship by default.
  \item[Part V, Operating Like an Expert] keys and models, workspaces and
    undo, reading the log, and using \nodwyn{} from anywhere.
  \item[The implementation map, Part~\ref{part:source-map}] guided routes
    through the Go host and every Scheme layer, followed by a worked
    explanation of this repository's own init file and development tools.
\end{description}

\section*{How to read it}

Every chapter has boxes:

\begin{tryit}
Something to type. Do it: the book is written to be tried, and every Scheme
program in it was run against the real interpreter before it was printed.
\end{tryit}

\begin{tip}
The kind of thing you would otherwise learn from a year of use.
\end{tip}

\begin{careful}
A sharp edge, and how not to cut yourself on it.
\end{careful}

\begin{puzzle}
A small problem to solve in Scheme. Answers are in
Appendix~\ref{app:answers}.
\end{puzzle}

When the book shows a session, lines starting with \code{you>} are what you
type; everything else is what \nodwyn{} prints. When it shows Scheme, the
result appears after a \code{;=>} comment:

\begin{scm}
(+ 1 2)          ;=> 3
\end{scm}

\section*{Where the truth lives}

This book was written from the repository: \file{docs/scheme.md} (the
language contract), \file{docs/manual.md}, the ADRs in \file{docs/adr/}, the
Scheme files in \file{internal/loop/}, and the source. When the book and the
repository disagree, the repository is right and the book has a bug. The
book-sized companion \file{nodwyn.pdf}, downloaded from GitHub Releases or
built locally at \file{docs/book/nodwyn.pdf}, tells the story of how
\nodwyn{} was designed and built; this one starts with using it well.
Both include the same implementation map so that either route can lead
to navigating and changing the codebase. The source stays in the repository;
the map points to the files and checks that establish the behaviour described.
